\section{Постановка и обозначения}
\label{sec:setup}

Спиновая частота частицы в магнитной структуре близка к $G\gamma$. Если равновесный уровень энергии различен для частиц сгустка, спиновые частоты расходятся, и поляризация теряет когерентность. В работе~\cite{Senichev2013} линейная по отклонению импульса часть этого эффекта подавлена высокочастотным полем, и остаётся второй порядок: квадратичная зависимость периода обращения от импульса и удлинение орбиты бетатронными колебаниями. Итог записан формулой (10) той работы. Цель этой записки --- вывести эту формулу из уравнений продольного движения, понять, откуда в ней берётся бетатронное слагаемое, и выяснить, что меняется в электростатическом кольце.

\subsection{Отклонение импульса и длина орбиты}

Отклонение импульса от синхронного
\begin{equation}
\label{eq:set-delta}
\delta=\frac{p-p_s}{p_s}.
\end{equation}
Индекс $s$ отмечает синхронную частицу: импульс $p_s$, скорость $v_s=\beta_s c$, лоренц-фактор $\gamma_s$, длина орбиты $C_s$, период $T_s=C_s/v_s$. Релятивистские $\beta_s$, $\gamma_s$ не смешиваются с функциями Твисса $\beta_{x,y}$, $\alpha_{x,y}$, $\gamma_{x,y}$, у которых всегда стоит индекс плоскости.

Длина замкнутой орбиты частицы с отклонением $\delta$ записывается рядом
\begin{equation}
\label{eq:set-C}
\frac{C}{C_s}=1+\alpha_0\delta+\alpha_1\delta^2+O(\delta^3).
\end{equation}
Коэффициенты $\alpha_0$, $\alpha_1$ --- коэффициенты ряда, а не производные: $\alpha_1=\tfrac12\,\mathrm{d}^2(C/C_s)/\mathrm{d}\delta^2$. Это определение работы~\cite{Senichev2013}, формула (15).

\subsection{Период и коэффициенты скольжения}

Относительное опоздание частицы за один оборот отсчитывается от периода синхронной частицы на том же обороте:
\begin{equation}
\label{eq:set-tau}
\tau\equiv\frac{T-T_s}{T_s}=\eta_0\delta+\eta_1\delta^2+\Lambda+\ldots
\end{equation}
Здесь $\eta_0$ и $\eta_1$ --- коэффициенты ряда по $\delta$, а $\Lambda$ --- вклад бетатронного движения, не зависящий от $\delta$ в удерживаемом порядке; он вводится в разделе~\ref{sec:beta}. Делитель в~\eqref{eq:set-tau} есть синхронный период $T_s$, как в определении (18.10) книги~\cite{Ng2006}. Почему именно он, разобрано в разделе~\ref{sec:ng}.

\subsection{Фаза и уравнения продольного движения}

Фаза $\varphi$ есть отставание частицы от синхронной, измеренное в радианах высокочастотного поля: она растёт, когда частица опаздывает. За один оборот
\begin{equation}
\label{eq:set-phase-map}
\Delta\varphi=\omega_\rf\,(T-T_s)=2\pi h\,\tau,
\end{equation}
где $\omega_\rf=2\pi h/T_s$, $h$ --- кратность. Напряжение зазора, которое видит частица, записываем как $V\sin(\varphi_s+\varphi)$. В накопителе без ускорения $\varphi_s=0$ или $\varphi_s=\pi$; устойчива та из двух точек, для которой
\begin{equation}
\label{eq:set-phis}
\cos\varphi_s=-\operatorname{sgn}\eta_0 .
\end{equation}
Ниже это проверяется (раздел~\ref{sec:mag}); такой выбор позволяет вести вывод одновременно для $\eta_0<0$ и $\eta_0>0$, что понадобится в электростатике.

За пролёт зазора полная энергия меняется на $eV\sin(\varphi_s+\varphi)=eV\cos\varphi_s\sin\varphi$. Деление на синхронный период $T_s=2\pi h/\omega_\rf$ (сглаживание) даёт систему
\begin{equation}
\label{eq:set-system}
\frac{\mathrm{d}\varphi}{\mathrm{d}t}=\omega_\rf\,\tau(\delta),
\qquad
\frac{\mathrm{d}\delta}{\mathrm{d}t}=K\cos\varphi_s\,\sin\varphi,
\qquad
K=\frac{eV\,\omega_\rf}{2\pi h\,\beta_s^2E_s},
\end{equation}
где использовано $\mathrm{d}E=\beta_s^2E_s\,\mathrm{d}\delta$ в синхронной точке. Время $t$ здесь --- часы синхронной частицы.

Знак фазы в~\cite{Senichev2013} противоположный: там $\mathrm{d}\varphi/\mathrm{d}t=-\omega_\rf\eta\delta$ при $\varphi_s=0$. Результат для равновесного сдвига от выбора знака не зависит (раздел~\ref{sec:mag}).

\subsection{Равновесный сдвиг}

Равновесным сдвигом $\delta_\eq$ называем среднее $\delta$ по синхротронному колебанию с амплитудой $\delta_m$ первой гармоники:
\begin{equation}
\label{eq:set-deq}
\delta_\eq=\langle\delta\rangle .
\end{equation}
Бетатронная фаза при этом уже усреднена: бетатронная частота много больше синхротронной.

\subsection{План}

Раздел~\ref{sec:mag} выводит $\tau(\delta)$ для магнитного кольца без бетатронного движения и сравнивает результат с формулами (1)--(10) работы~\cite{Senichev2013}. Раздел~\ref{sec:sol} решает продольные уравнения методом Боголюбова--Митропольского и оценивает поправки. Раздел~\ref{sec:beta} объясняет, как в период входит бетатронное удлинение орбиты. Раздел~\ref{sec:ng} разбирает определения коэффициента скольжения. Раздел~\ref{sec:es} переносит вывод на электростатическое кольцо, раздел~\ref{sec:f25} --- сравнение с формулой (25) работы~\cite{Senichev2013}. Схема метода Боголюбова--Митропольского вынесена в приложение~\ref{app:bm}.
